% from ega0-22-fr.tex
\oldpage[0\textsubscript{IV-1}]{212}

est injectif. Considérons encore le diagramme commutatif
\begin{equation}
\label{0.22.7.3.2-fr}
\xymatrix{
  (\mathfrak q/\mathfrak q^2)\otimes_AL
    \ar[r]^{j_0} \ar[d]_\alpha &
  \Omega^1_{A/k}\otimes_AL \ar[r] \ar@{=}[d] &
  \widehat\Omega^1_{A/k'}\otimes_AL \ar@{=}[d] \\
  (\mathfrak p/\mathfrak p^2)\otimes_AL
    \ar[r]_{i_0} &
  \Omega^1_{A/k}\otimes_AL \ar[r] &
  \widehat\Omega^1_{A/k'}\otimes_AL
}
\tag{22.7.3.2}
\end{equation}

Comme $j_0$ est injectif, on a $\operatorname{Ker}(i_0\circ\alpha)=0$. Mais comme
$A_{\mathfrak p}$ est un anneau régulier et que l'hypothèse a) implique que
$B_{\mathfrak p}=A_{\mathfrak p}/\mathfrak qA_{\mathfrak p}$ est aussi un
anneau régulier (19.6.5), le lemme (22.7.2.4) montre que $\alpha$ est injectif,
d'où $\alpha^{-1}(\operatorname{Ker} i_0)=0$. Si $i'$ est le composé des homomorphismes de la
seconde ligne de (22.7.3.2), on a de même
$\operatorname{Ker}(j')=\alpha^{-1}(\operatorname{Ker}(i'))$, et tout revient donc à montrer que, pour un
choix convenable de $k'$, on a
\begin{equation}
\label{0.22.7.3.3-fr}
  \operatorname{Ker}(i')=\operatorname{Ker}(i).
  \tag{22.7.3.3}
\end{equation}

Mais le même raisonnement que dans (22.7.2) montre que l'on a une suite exacte
\[
  (\mathfrak p/\mathfrak p^2)\otimes_AL\xrightarrow{i'}
  \widehat\Omega^1_{A/k'}\otimes_AL\longrightarrow
  \widehat\Omega^1_{(A/\mathfrak p)/k'}\otimes_{A/\mathfrak p}L
  \longrightarrow0.
\]
Or, les hypothèses faites permettent d'appliquer (21.9.8), donc il existe une
sous-extension $k'$ de $K_0$ contenant $K_0^p$, telle que
$[K_0:k']<+\infty$, et pour laquelle on a
\[
  \operatorname{rg}_L\bigl(\widehat\Omega^1_{(A/\mathfrak p)/k'}
  \otimes_{A/\mathfrak p}L\bigr)
  =\dim(A/\mathfrak p)+\operatorname{rg}_L\Upsilon_{L/k}
  +\operatorname{rg}_{K_0}\Omega^1_{K_0/k'}.
\]
Mais on a
\[
  \dim(A/\mathfrak p)=\dim(A)-\dim(A_{\mathfrak p})
\]
par (16.5.11),
\[
  \dim(A_{\mathfrak p})
  =\operatorname{rg}_L\bigl((\mathfrak p/\mathfrak p^2)\otimes_AL\bigr)
\]
puisque $A_{\mathfrak p}$ est régulier (17.1.1), et enfin
\[
  \dim(A)+\operatorname{rg}_{K_0}\Omega^1_{K_0/k'}
  =\operatorname{rg}_L\bigl(\widehat\Omega^1_{A/k'}\otimes_AL\bigr)
\]
en vertu de (21.9.2)~; on obtient donc
\begin{equation}
\label{0.22.7.3.4-fr}
  \operatorname{rg}_L(\operatorname{Ker}(i'))=\operatorname{rg}_L\Upsilon_{L/k}.
  \tag{22.7.3.4}
\end{equation}

D'autre part, le corps résiduel de $A_{\mathfrak p}$ étant formellement lisse
sur le corps premier de $k$, on a une suite exacte (20.6.22.1)
\[
  \Upsilon_{L/k}\xrightarrow{\chi_{A_{\mathfrak p}/k}}
  (\mathfrak pA_{\mathfrak p})/(\mathfrak pA_{\mathfrak p})^2
  \xrightarrow{i}\Omega^1_{A_{\mathfrak p}/k}
  \otimes_{A_{\mathfrak p}}L,
\]
et comme $A_{\mathfrak p}$ est une $k$-algèbre formellement lisse pour la
topologie $\mathfrak p$-préadique, il suit de (22.2.2) que
$\chi_{A_{\mathfrak p}/k}$ est un homomorphisme injectif, donc $\operatorname{Ker}(i)$ est
isomorphe à $\Upsilon_{L/k}$, et compte tenu de (22.7.3.4) et de ce que
$\operatorname{Ker}(i)\subset\operatorname{Ker}(i')$, cela achève de prouver (22.7.3.3) et par suite le
théorème.
\end{proof}

